% 2-preliminarii.tex starts here
\chapter{Preliminarii}\label{ch:prelim}

\section*{Contextul tehnologic}
Proiectul \textit{Kookio} se bazează pe arhitectura unitară \textbf{PAN} (Prisma–Analog–Nx), care integrează fluxul complet de dezvoltare \texttt{TypeScript} de la interfață până la bazele de date.
Această arhitectură facilitează modularitatea, testarea automată și extinderea sistemului, fiind susținută de instrumente moderne pentru orchestrare, randare hibridă și persistență tipizată \cite{nx,analog,prisma}.
O prezentare detaliată a manifestului, a avantajelor comparative și a direcțiilor de evoluție este centralizată în \textbf{Anexa~\ref{ch:appendix-manifest}}.

\section{Preliminarii}

\subsection{Arhitecturi web moderne}

Aplicațiile web moderne au evoluat semnificativ față de modelul clasic client--server, în care browserul funcționa exclusiv ca un consumator de API-uri REST.
Cerințele actuale privind performanța, experiența utilizatorului și optimizarea pentru motoarele de căutare au condus la apariția unor arhitecturi hibride, care combină randarea pe client cu randarea pe server \cite{analog,nitro}.

Printre aceste abordări se numără:
\begin{itemize}
  \item \textbf{Client-Side Rendering (CSR)}, unde aplicația este randată integral în browser;
  \item \textbf{Server-Side Rendering (SSR)}, unde conținutul HTML este generat pe server pentru fiecare cerere;
  \item \textbf{Static Site Generation (SSG)}, unde paginile sunt generate la build-time;
  \item \textbf{arhitecturi hibride}, care combină aceste tehnici în funcție de tipul conținutului.
\end{itemize}

În contextul proiectului Kookio, este adoptată o arhitectură full-stack modernă, care integrează frontend-ul și backend-ul într-un flux unitar de dezvoltare, permițând partajarea tipurilor și a logicii de validare între straturi \cite{trpc}.
Această abordare reduce complexitatea integrării și crește coerența aplicației.

\subsection{Angular și randare server-side}

Angular este un framework matur pentru dezvoltarea aplicațiilor web, bazat pe TypeScript și orientat către arhitecturi scalabile \cite{angular_signals}.
În mod tradițional, Angular a fost utilizat predominant în scenarii de tip client-side rendering, însă versiunile recente extind semnificativ suportul pentru randare server-side.

În cadrul proiectului Kookio, Angular este utilizat împreună cu Analog, un meta-framework care adaugă suport nativ pentru:
\begin{itemize}
  \item Server-Side Rendering (SSR),
  \item rutare bazată pe fișiere,
  \item integrarea logicii de server în aceeași aplicație \cite{analog}.
\end{itemize}

Randarea server-side permite generarea conținutului HTML direct pe server, îmbunătățind timpul de afișare inițială și compatibilitatea cu motoarele de căutare \cite{nitro}.
Această caracteristică este esențială pentru aplicații orientate către conținut, precum Kookio, unde lizibilitatea inițială și performanța percepută sunt importante.

Analog facilitează această abordare fără a separa artificial frontend-ul și backend-ul, păstrând o bază de cod unificată și coerentă.

\subsection{Monorepo și orchestrare cu Nx}

Proiectul este organizat sub forma unui monorepo, utilizând instrumentul Nx.
Abordarea monorepo permite gestionarea mai multor aplicații și biblioteci într-un singur repository, cu reguli explicite de dependență și caching la nivel de task \cite{nx}.

În cadrul proiectului Kookio, această structură facilitează separarea clară între:
\begin{itemize}
  \item aplicația web (frontend și server-side logic),
  \item stratul de acces la date (Prisma Client),
  \item validarea și contractele de date (Zod),
  \item utilitare și date mock partajate.
\end{itemize}

Nx oferă avantaje importante precum:
\begin{itemize}
  \item execuția incrementală a task-urilor (build, test, lint),
  \item cache local și determinism al build-urilor,
  \item izolare logică între proiecte prin reguli de boundary \cite{nx}.
\end{itemize}

Această abordare este deosebit de potrivită pentru aplicații full-stack moderne, unde frontend-ul și backend-ul evoluează în paralel și partajează tipuri și contracte comune.

\subsection{tRPC și partajarea tipurilor}

Pentru comunicarea între client și server, proiectul utilizează tRPC, un framework care permite definirea endpoint-urilor API într-un mod complet tipat, fără a necesita generarea manuală de clienți sau specificații intermediare \cite{trpc}.

În această arhitectură, procedurile tRPC sunt definite pe server și consumate direct în aplicația Angular, beneficiind de:
\begin{itemize}
  \item inferență completă a tipurilor TypeScript,
  \item validare automată a input-urilor prin Zod \cite{zod},
  \item eliminarea discrepanțelor dintre contractele client și server.
\end{itemize}

Această abordare reduce semnificativ erorile de integrare și accelerează dezvoltarea, fiind adecvată în special aplicațiilor full-stack TypeScript.

\subsection{Persistență: Prisma și CockroachDB}

Stratul de persistență este construit folosind Prisma ca ORM și CockroachDB ca sistem de gestiune a bazei de date \cite{prisma,cockroachdb}.
CockroachDB este un sistem NewSQL distribuit, compatibil cu protocolul PostgreSQL, care oferă consistență tranzacțională și scalabilitate orizontală \cite{postgresql}.

Prisma este utilizat pentru:
\begin{itemize}
  \item definirea schemei bazei de date într-un limbaj declarativ,
  \item generarea automată a clientului TypeScript,
  \item gestionarea migrațiilor între versiuni ale schemei.
\end{itemize}

Alegerea acestei combinații permite proiectului să rămână compatibil cu un ecosistem SQL matur, dar și pregătit pentru scenarii de scalare viitoare, fără modificări semnificative ale codului aplicației.

\subsection{Testare automată: Vitest și Playwright}

Testarea automată reprezintă un element esențial în dezvoltarea aplicațiilor moderne, asigurând corectitudinea funcțională și stabilitatea în timp.
În cadrul proiectului Kookio sunt utilizate două tipuri complementare de testare: teste unitare și teste end-to-end.

Pentru testarea unitară și de integrare este utilizat Vitest, un framework de testare rapid, compatibil cu ecosistemul Vite și TypeScript \cite{vitest,vite}.
Vitest permite testarea componentelor Angular, a utilitarelor și a logicii de business într-un mediu izolat, cu suport pentru mocking și rulare incrementală.

Testarea end-to-end este realizată cu ajutorul Playwright, care permite simularea interacțiunilor reale ale utilizatorului în browser \cite{playwright}.
Acest tip de testare validează fluxurile critice ale aplicației, precum încărcarea paginilor, interacțiunea cu formularul și comunicarea cu API-ul server-side.

Combinarea celor două abordări oferă o acoperire echilibrată:
\begin{itemize}
  \item Vitest validează corectitudinea logicii interne,
  \item Playwright verifică comportamentul aplicației din perspectiva utilizatorului final.
\end{itemize}

Această strategie contribuie la creșterea încrederii în stabilitatea aplicației și facilitează evoluția incrementală a codului.

\subsection{Managementul configurației și mediilor}

Variabilele de mediu sunt gestionate centralizat în platforma de deployment Vercel, care acționează ca sursă unică de adevăr pentru configurația aplicației.
Această abordare permite separarea clară a mediilor (dezvoltare, testare, producție) și reduce riscul expunerii accidentale a secretelor \cite{vercel_cli}.

Pentru rularea locală, Vercel CLI este utilizat pentru a sincroniza temporar variabilele de mediu active prin comanda \texttt{vercel env pull}.
Astfel, mediul local rămâne aliniat cu mediul de deploy, fără a necesita fișiere \texttt{.env} persistente în repository.
% 2-preliminarii.tex ends here